\subsection{Entrelazamiento APP--Witt, vecino de Leech y orbifold de orden \(3\)}

Las doce fibras llevan las etiquetas
\[
 (i,\mu,\varepsilon)\in
 \mathbb Z/3\mathbb Z\times\{0,1\}\times\{0,1\},
 \qquad
 \lambda(i,\mu,\varepsilon)=4i+2\mu+\varepsilon\pmod{12}.
\]
La aplicación \(\lambda\) es biyectiva y, fijados trivializaciones y marcos,
los entrelazadores satisfacen
\[
 g_{\lambda(i,\mu,\varepsilon)}T_{i\mu\varepsilon}
 =T_{i\mu\varepsilon}C^{(-1)^\mu}.
\]
Su suma directa es un isomorfismo \(C_3\)-equivariante
\[
 T_\lambda:W_{\mathrm{APP}}\otimes\mathbb C
 \xrightarrow{\ \cong\ }
 \bigoplus_{b=0}^{11}(A_{2,b}\otimes\mathbb C).
\]

Sea \(N=N(A_2^{12})\), sea \(v=v_\alpha\) el vector radial de norma \(54\),
y definamos
\[
 N_0=\ker\bigl(\langle\,\cdot\,,v\rangle\bmod3\bigr),
 \qquad N_v=N_0+\mathbb Z\frac v3.
\]
El cálculo reticular verifica que \(N_v\cong\Lambda_{24}\), que la
isometría \(g_c\) tiene orden tres y no posee puntos fijos no nulos, y que
\[
 \frac{g_cv-v}{3},\quad\frac{g_c^{-1}v-v}{3}\in N_0.
\]
Por ello \(g_cN_v=N_v\). Su elevación estándar al VOA reticular tiene peso
torcido \(4/3\) y tipo cero. Los teoremas clásicos de orbifold cíclico dan
\[
 V_{(3)}:=\operatorname{Orb}_{\mathbb Z/3}
       (V_{N_v},\widehat g_c)\cong V^\natural,
 \qquad T_{3B}=t_3+12.
\]
El carácter se recupera por el promedio de los nueve sectores y por la
uniformización de nivel tres:
\[
 \operatorname{ch}_{V_{(3)}}
 =\frac13\sum_{i,j\in\mathbb Z/3}Z_{i,j}=J,
\]
\[
 J=t_3+12+\frac{10\,3^9}{t_3}
       +\frac{4\,3^{14}}{t_3^2}+\frac{3^{18}}{t_3^3}.
\]
En particular,
\[
 [q]J=54+10\cdot3^9=196884=54(5\cdot729+1).
\]

\subsection{Construcciones de órdenes \(2\) y \(3\) del álgebra de operadores de vértice \texorpdfstring{\(V^\natural\)}{V natural}}

Fijemos la extensión central y el cociclo de Leech
\[
 1\longrightarrow\{\pm1\}\longrightarrow\widehat\Lambda_{24}
 \longrightarrow\Lambda_{24}\longrightarrow1,
 \qquad
 V_\Lambda=M(1)\otimes\mathbb C_\varepsilon[\Lambda_{24}].
\]
La negación \(\lambda\mapsto-\lambda\) se eleva a \(\theta\), y
\[
 t_2(\tau)=\operatorname{Tr}_{V_\Lambda}
  \left(\theta q^{L_0-1}\right)
 =q^{-1}\prod_{n\ge1}(1+q^n)^{-24}.
\]
La presentación FLM es
\[
 V_{(2)}:=(V_\Lambda)^+\oplus(V_\Lambda^T)^+\cong V^\natural.
\]

\begin{theorem}[Equivalencia tipada de las presentaciones de órdenes \(3\) y \(2\)]
\label{thm:cierre-vtm-equivalencia-voa-23}
Sean
\[
 \phi_3:V_{(3)}\xrightarrow{\ \cong\ }V^\natural,
 \qquad
 \phi_2:V_{(2)}\xrightarrow{\ \cong\ }V^\natural
\]
los isomorfismos graduados proporcionados por los teoremas de orbifold
aplicados a los datos anteriores. Entonces
\[
 \boxed{\Phi_{23}:=\phi_2^{-1}\circ\phi_3:
 V_{(3)}\xrightarrow{\ \cong\ }V_{(2)}}
\]
es un isomorfismo graduado de álgebras de operadores de vértice. Además,
\[
 a\longmapsto\Phi_{23}a\Phi_{23}^{-1}
\]
identifica sus grupos de automorfismos y conserva las trazas graduadas.
\end{theorem}

\begin{proof}
La composición de isomorfismos de VOA es un isomorfismo de VOA y preserva el
vacío, el vector conforme, los operadores de vértice y la gradación por
\(L_0\). La conjugación transporta automorfismos. Para cada automorfismo
\(a\), la invariancia de la traza bajo conjugación da la igualdad de las
series graduadas correspondientes.
\end{proof}

El isomorfismo \(\Phi_{23}\) depende de las identificaciones \(\phi_2\) y
\(\phi_3\); no se afirma canonicidad absoluta. La construcción de orden dos
conserva además el contenido exclusivo
\[
 T_{2A}=t_2+24+\frac{2^{12}}{t_2},
 \qquad
 \boxed{j=\frac{(t_2+256)^3}{t_2^2}},
\]
que completa la uniformización modular de orden dos.

Una vez construido \(V^\natural\), las tres publicaciones se enuncian sin
serializarlas indebidamente:
\[
 \operatorname{Aut}(V^\natural)\cong\mathbb M,
 \qquad
 \operatorname{ch}_{V^\natural}=J,
 \qquad
 g\longmapsto T_g(\tau)
 =\operatorname{Tr}_{V^\natural}\!\left(gq^{L_0-1}\right).
\]
El teorema de Moonshine coordina la acción graduada con las funciones de
género cero y sus identidades de replicabilidad. La acción del Monstruo vive
en \(V^\natural\); que no sea una permutación puntual de cifras es un control
de tipos final, no la tesis del corredor.

